% Generated draft: requires mathematical and source review.
\documentclass[11pt]{article}
\usepackage[margin=25mm]{geometry}
\usepackage{fontspec}
\setmainfont{Latin Modern Roman}
\usepackage{amsmath,amssymb,mathtools,unicode-math}
\setmathfont{Latin Modern Math}
\usepackage{enumitem,needspace,xurl,hyperref}
\setlength{\parindent}{0pt}
\setlength{\parskip}{5pt}
\setlength{\emergencystretch}{3em}
\newcommand{\N}{\mathbb{N}}
\newcommand{\Z}{\mathbb{Z}}
\newcommand{\Q}{\mathbb{Q}}
\newcommand{\R}{\mathbb{R}}
\newcommand{\C}{\mathbb{C}}
\newcommand{\F}{\mathbb{F}}
\newcommand{\A}{\mathbb{A}}
\newcommand{\PP}{\mathbb{P}}
\newcommand{\E}{\mathbb{E}}
\newcommand{\eps}{\varepsilon}
\newcommand{\dd}{\,\mathrm{d}}
\begin{document}
\section*{UnsolvedMath ID: 800205}
\textbf{Source problem code:} catalogue entry
\subsection{Definitions and mathematical statement}
Let \(V\) be the germ at \(0\) of an irreducible complex variety in \(\C^n\). For a sufficiently small \(\epsilon>0\), choose a representative of this germ on \(B_\epsilon:=\{z\in\C^n:|z|<\epsilon\}\), and put
\[
E_\epsilon:=V\cap\R^n\cap B_\epsilon.
\]
Here \(|\cdot|\) denotes the Euclidean norm, and \(\Im z\in\R^n\) is the vector of imaginary parts of the coordinates of \(z\). The supplied text defines
\[
U_{E_\epsilon}(z):=\sup\bigl\{u(z):u\leq0\text{ on }E_\epsilon,\quad 0\leq u\leq1\text{ on }B_\epsilon\bigr\},\qquad z\in B_\epsilon.
\]
The source does not specify the function class over which this supremum is taken. If its displayed conditions are read literally, the admissible functions are all maps \(u:B_\epsilon\to\R\) satisfying those conditions; no continuity, plurisubharmonicity, or other regularity condition is stated.

For a fixed sufficiently small \(\epsilon\), the displayed classification target is to characterize those germs \(V\) for which
\[
\exists A>0\quad\forall z\in B_{\epsilon/2},\qquad U_{E_\epsilon}(z)\leq A|\Im z|.\tag{\(\alpha\)}
\]
The displayed domain of this estimate is the ambient ball \(B_{\epsilon/2}\). The source does not state whether \(A\) may depend on \(\epsilon\), or whether a single \(A\) must work for every sufficiently small \(\epsilon\). Consequently, the intended classification remains incompletely specified until the admissible function class and the quantification of \(\epsilon\) and \(A\) are clarified. The source also refers to Section 7.2, whose contents are not supplied.
\subsection{Short English statement}
Classify irreducible complex variety germs whose real locus has the stated extremal-function bound by the distance to \(\R^n\). The supplied formulation omits the admissible function class and leaves the dependence of the bound on \(\epsilon\) unresolved.
\subsection{Sources}
\begin{itemize}
\item Local UnsolvedMath registry: permanent ID \(800205\), source code \texttt{catalogue entry}, title \emph{Several Complex Variables Problem List --- Problem 1}. The supplied record attributes the problem to Buzzard, Fornaess, and Krantz, with proposed year 1995, and labels it open; this status is supplied metadata, not an independently verified assessment.
\item Buzzard, Fornaess, Krantz, \emph{The Several Complex Variables Problem List} (1995), Problem 1. Supplied source URL: \url{https://arxiv.org/abs/math/9509202}. The excerpt marks the item ``New'' and refers to Section 7.2.
\item Supplied provenance records access on \texttt{2026-10-16} and extraction method \texttt{arxiv\_latex}. This is source-reported provenance; no independent source visit or verification is claimed.
\end{itemize}
\end{document}
